% ch12_monoidal_enriched.tex

\chapter{Monoidal and Enriched Categories}

\begin{overviewbox}
The categories you have encountered so far have \emph{sets} of morphisms between
any two objects. The hom-set $\mathcal{C}(a, b)$ is just a set: a collection of
arrows from $a$ to $b$, with no additional structure.

But in many important categories, those hom-sets carry extra structure. In the
category of abelian groups, $\operatorname{Hom}(A, B)$ is itself an abelian group.
In a category of chain complexes, morphism spaces are chain complexes. In a
2-category, the morphisms between two objects form a whole category.

The unifying concept is \textbf{enriched category theory}: a category enriched
over a monoidal category $\mathcal{V}$ replaces its hom-sets with objects of
$\mathcal{V}$. To make this work, $\mathcal{V}$ must itself be a
\textbf{monoidal category} — a category equipped with a tensor product that
plays the role of the cartesian product for sets.

This chapter builds monoidal categories from scratch, develops the coherence
theorem that makes them well-behaved, introduces symmetric and closed monoidal
categories, and then constructs the general theory of $\mathcal{V}$-enriched
categories.
\end{overviewbox}

% ─────────────────────────────────────────────────────────────────────────────
\section{Monoidal Categories}

\subsection{Informally}

A \term{monoidal category} is a category with a way to ``multiply'' objects
and morphisms together. The archetype is $(\mathbf{Set}, \times, \{*\})$: sets
with cartesian product, where the one-element set $\{*\}$ acts as the identity
for $\times$. But many other examples exist:

\begin{itemize}
  \item $(\mathbf{Ab}, \otimes_{\mathbb{Z}}, \mathbb{Z})$: abelian groups with tensor product.
  \item $(\mathbf{Vect}_k, \otimes_k, k)$: vector spaces over a field $k$.
  \item $(\mathbf{Set}, +, \emptyset)$: sets with disjoint union.
  \item $([\mathcal{C}, \mathcal{C}], \circ, \operatorname{Id})$: endofunctors on a
    category, with composition.
\end{itemize}

The defining data is: a category $\mathcal{V}$, a functor
$\otimes : \mathcal{V} \times \mathcal{V} \to \mathcal{V}$ (the tensor product),
and a unit object $I \in \mathcal{V}$, together with natural isomorphisms
expressing associativity and left/right unitality.

\subsection{Expanding}

The tricky part is that associativity and unitality do not need to be equalities —
they can be natural isomorphisms (``coherent up to isomorphism''):
\begin{align*}
  \alpha_{A,B,C} &: (A \otimes B) \otimes C \xrightarrow{\sim} A \otimes (B \otimes C)
    \qquad \text{(associator)}, \\
  \lambda_A &: I \otimes A \xrightarrow{\sim} A
    \qquad \text{(left unitor)}, \\
  \rho_A &: A \otimes I \xrightarrow{\sim} A
    \qquad \text{(right unitor)}.
\end{align*}

Having these isomorphisms opens a question: if there are many ways to parenthesize
an expression like $A \otimes B \otimes C \otimes D$, do all the paths between
any two parenthesizations agree? \textbf{Mac Lane's coherence theorem} says yes:
any two composites of associators and unitors between the same parenthesizations
are equal. This means we can compute in monoidal categories as if the tensor
product were strictly associative — writing $A \otimes B \otimes C$ without
parentheses — as long as we are careful about what ``equal'' means.

\subsection{Formalizing}

\begin{definition}
A \term{monoidal category} $(\mathcal{V}, \otimes, I, \alpha, \lambda, \rho)$
consists of:
\begin{itemize}
  \item A category $\mathcal{V}$.
  \item A \term{tensor product} functor $\otimes : \mathcal{V} \times \mathcal{V} \to \mathcal{V}$.
  \item A \term{unit object} $I \in \mathcal{V}$.
  \item Natural isomorphisms:
  \begin{align*}
    \alpha_{A,B,C} &: (A \otimes B) \otimes C \xrightarrow{\sim} A \otimes (B \otimes C),\\
    \lambda_A &: I \otimes A \xrightarrow{\sim} A, \\
    \rho_A   &: A \otimes I \xrightarrow{\sim} A.
  \end{align*}
\end{itemize}
These must satisfy the \term{pentagon axiom}:
\begin{center}
\begin{tikzcd}[column sep=small]
& (A \otimes B) \otimes (C \otimes D)
    \arrow[dr, "\alpha_{A,B,C\otimes D}"]
& \\
((A \otimes B) \otimes C) \otimes D
    \arrow[ur, "\alpha_{A\otimes B,C,D}"]
    \arrow[d, "\alpha_{A,B,C} \otimes \operatorname{id}_D"']
& & A \otimes (B \otimes (C \otimes D)) \\
(A \otimes (B \otimes C)) \otimes D
    \arrow[rr, "\alpha_{A,B\otimes C,D}"']
& & A \otimes ((B \otimes C) \otimes D)
    \arrow[u, "\operatorname{id}_A \otimes \alpha_{B,C,D}"']
\end{tikzcd}
\end{center}
and the \term{triangle axiom}:
\begin{center}
\begin{tikzcd}
(A \otimes I) \otimes B
    \arrow[rr, "\alpha_{A,I,B}"]
    \arrow[dr, "\rho_A \otimes \operatorname{id}_B"']
& & A \otimes (I \otimes B)
    \arrow[dl, "\operatorname{id}_A \otimes \lambda_B"]
\\
& A \otimes B &
\end{tikzcd}
\end{center}
\end{definition}

\begin{howtosay}
The tensor product $A \otimes B$ is read ``$A$ tensor $B$.'' The unit $I$ may be
called the \emph{monoidal unit} or \emph{unit object}. A monoidal category where
$\alpha$, $\lambda$, $\rho$ are all identity morphisms is called \term{strict}.
\end{howtosay}

\begin{theorem}[Mac Lane's Coherence Theorem]
Every monoidal category is monoidally equivalent to a strict monoidal category.
Moreover, every diagram built from instances of $\alpha$, $\lambda$, $\rho$
and their inverses that connects the same two parenthesizations commutes.
\end{theorem}

Coherence allows us to write $A \otimes B \otimes C$ without parentheses and
to move freely between parenthesizations without tracking the specific isomorphisms.

% ─────────────────────────────────────────────────────────────────────────────
\section{Symmetric and Closed Monoidal Categories}

\subsection{Formalizing}

Two additional structures enrich monoidal categories significantly.

\begin{definition}
A monoidal category $(\mathcal{V}, \otimes, I)$ is \term{symmetric} if it has
a natural isomorphism
\[
  \sigma_{A,B} : A \otimes B \xrightarrow{\sim} B \otimes A
\]
satisfying: $\sigma_{B,A} \circ \sigma_{A,B} = \operatorname{id}_{A \otimes B}$,
and compatibility with $\alpha$ expressed by the hexagon axioms.
\end{definition}

\begin{definition}
A monoidal category $(\mathcal{V}, \otimes, I)$ is \term{closed} if for each
object $B$, the functor $(-) \otimes B : \mathcal{V} \to \mathcal{V}$ has a right
adjoint, written $[B, -]$ or $B \multimap (-)$\sayit{$B$ hom dash}:
\[
  \mathcal{V}(A \otimes B,\, C) \;\cong\; \mathcal{V}(A,\, [B, C]).
\]
The object $[B, C]$ is called the \term{internal hom} of $B$ and $C$.
\end{definition}

\begin{howtosay}
$[B, C]$ is read ``$B$ to $C$'' or ``the internal hom from $B$ to $C$.'' In
$\mathbf{Set}$, $[B, C] = C^B$ is the set of functions from $B$ to $C$; the
adjunction becomes the familiar currying bijection
$\operatorname{Hom}(A \times B, C) \cong \operatorname{Hom}(A, C^B)$.
\end{howtosay}

The category $\mathbf{Set}$ with cartesian product is symmetric monoidal closed.
The category $\mathbf{Ab}$ with $\otimes_{\mathbb{Z}}$ is symmetric monoidal closed,
with internal hom $[A, B] = \operatorname{Hom}_{\mathbb{Z}}(A, B)$.

% ─────────────────────────────────────────────────────────────────────────────
\section{Enriched Categories}

\subsection{Informally}

Fix a monoidal category $(\mathcal{V}, \otimes, I)$. A \term{$\mathcal{V}$-enriched
category} $\mathcal{C}$, or \term{$\mathcal{V}$-category}, is like an ordinary
category except that its hom-sets are replaced by objects of $\mathcal{V}$.

The idea is simple: instead of a \emph{set} $\mathcal{C}(a, b)$ of morphisms
from $a$ to $b$, we have an \emph{object} $\mathcal{C}(a, b) \in \mathcal{V}$.
Composition becomes a morphism in $\mathcal{V}$:
\[
  \circ_{a,b,c} : \mathcal{C}(b, c) \otimes \mathcal{C}(a, b) \;\to\; \mathcal{C}(a, c).
\]
Identity becomes a morphism from the unit:
\[
  j_a : I \;\to\; \mathcal{C}(a, a).
\]

\subsection{Expanding}

Here are the key examples:

\begin{itemize}
  \item \textbf{$\mathbf{Set}$-enriched categories} are ordinary categories.
  \item \textbf{$\mathbf{Ab}$-enriched categories} are \term{preadditive categories}:
    the hom-sets are abelian groups and composition is bilinear.
  \item \textbf{$\mathbf{Vect}_k$-enriched categories} are \term{$k$-linear categories}.
  \item \textbf{$\mathbf{Cat}$-enriched categories} are \term{2-categories}:
    hom-objects are themselves categories (categories of morphisms between two objects).
  \item \textbf{$([0,\infty], \geq, +, 0)$-enriched categories} are \term{Lawvere metric spaces}:
    objects are points, and $\mathcal{C}(a,b)$ is the ``distance from $a$ to $b$.''
\end{itemize}

The last example is striking. A monoidal category need not be cartesian. The
poset $([0,\infty], \geq)$ — real numbers ordered by $\geq$, with tensor product
$a + b$ and unit $0$ — is a monoidal category. An enriched category over it is
a set with a distance function satisfying:
\[
  0 \geq d(a, a), \qquad d(b,c) + d(a,b) \geq d(a,c).
\]
That is: identity has distance $0$, and the triangle inequality. This is a
(generalized) metric space.

\subsection{Formalizing}

\begin{definition}
Let $(\mathcal{V}, \otimes, I)$ be a monoidal category. A
\term{$\mathcal{V}$-enriched category} (or \term{$\mathcal{V}$-category}) $\mathcal{C}$
consists of:
\begin{itemize}
  \item A collection of \term{objects} $\operatorname{ob}(\mathcal{C})$.
  \item For each pair $a, b \in \operatorname{ob}(\mathcal{C})$, an object
    $\mathcal{C}(a,b) \in \mathcal{V}$, called the \term{hom-object}.
  \item For each triple $a, b, c$, a \term{composition morphism}
    \[
      \circ_{a,b,c} : \mathcal{C}(b,c) \otimes \mathcal{C}(a,b) \;\to\; \mathcal{C}(a,c).
    \]
  \item For each object $a$, an \term{identity morphism}
    \[
      j_a : I \;\to\; \mathcal{C}(a,a).
    \]
\end{itemize}
These must satisfy associativity and unit laws expressed as commutative diagrams
in $\mathcal{V}$, analogous to the ordinary category axioms.
\end{definition}

% ─────────────────────────────────────────────────────────────────────────────
\section{$\mathcal{V}$-Functors and $\mathcal{V}$-Natural Transformations}

\subsection{Formalizing}

\begin{definition}
A \term{$\mathcal{V}$-functor} $F : \mathcal{C} \to \mathcal{D}$ between
$\mathcal{V}$-categories consists of:
\begin{itemize}
  \item A function $F : \operatorname{ob}(\mathcal{C}) \to \operatorname{ob}(\mathcal{D})$.
  \item For each pair $a, b \in \mathcal{C}$, a morphism in $\mathcal{V}$:
    \[
      F_{a,b} : \mathcal{C}(a, b) \;\to\; \mathcal{D}(Fa, Fb).
    \]
\end{itemize}
These must be compatible with composition and identities via commutative diagrams
in $\mathcal{V}$.
\end{definition}

\begin{definition}
A \term{$\mathcal{V}$-natural transformation} $\alpha : F \Rightarrow G$ between
$\mathcal{V}$-functors $F, G : \mathcal{C} \to \mathcal{D}$ assigns to each
$a \in \mathcal{C}$ a morphism $\alpha_a : I \to \mathcal{D}(Fa, Ga)$ in $\mathcal{V}$,
satisfying the naturality condition expressed as a commutative square in $\mathcal{V}$.
\end{definition}

\begin{howtosay}
A $\mathcal{V}$-natural transformation has its ``components'' as morphisms out of
the unit $I$ into hom-objects, rather than as elements of a set. In the
$\mathbf{Set}$-enriched case, $I = \{*\}$ and a morphism $\{*\} \to \mathcal{D}(Fa, Ga)$
is the same as picking an element of $\mathcal{D}(Fa, Ga)$, recovering the ordinary
definition.
\end{howtosay}

\subsection{Self-Enrichment}

A monoidal category $\mathcal{V}$ that is also closed can be enriched over itself.
The hom-object between two objects $a, b$ is the internal hom $[a, b] \in \mathcal{V}$,
and composition is given by the morphism:
\[
  [b, c] \otimes [a, b] \;\to\; [a, c]
\]
obtained via the adjunction for closed structure. This self-enrichment is what
makes closed monoidal categories the natural foundation for enriched category theory.

\subsection{Reorganizing: Ordinary CT as Enriched CT}

Every concept from Volumes I can be enriched:

\begin{center}
\renewcommand{\arraystretch}{1.4}
\begin{tabular}{ll}
\hline
\textbf{Ordinary} & \textbf{Enriched over $\mathcal{V}$} \\
\hline
Category & $\mathcal{V}$-category \\
Functor & $\mathcal{V}$-functor \\
Natural transformation & $\mathcal{V}$-natural transformation \\
Adjunction & $\mathcal{V}$-adjunction \\
Limit / Colimit & $\mathcal{V}$-limit / $\mathcal{V}$-colimit (weighted) \\
Yoneda lemma & $\mathcal{V}$-Yoneda lemma \\
Kan extension & $\mathcal{V}$-Kan extension \\
\hline
\end{tabular}
\end{center}

The passage from ordinary to enriched is not merely generalization for its own
sake. Many important categories are naturally enriched (abelian groups,
chain complexes, spectra in homotopy theory), and their enriched structure
carries essential information that the underlying ordinary structure loses.

% ─────────────────────────────────────────────────────────────────────────────
\section{Exercises}

\begin{exercisesbox}
\begin{qa}
\qaitem{What is a monoidal category?}
       {A category $\mathcal{V}$ with a tensor product functor $\otimes$, a unit
        object $I$, and natural isomorphisms $\alpha$ (associator), $\lambda$
        (left unitor), $\rho$ (right unitor), satisfying pentagon and triangle axioms.}

\qaitem{What are the three natural isomorphisms in a monoidal category?}
       {Associator $\alpha_{A,B,C} : (A\otimes B)\otimes C \cong A\otimes(B\otimes C)$;
        left unitor $\lambda_A : I\otimes A \cong A$; right unitor $\rho_A : A\otimes I \cong A$.}

\qaitem{What does the pentagon axiom say?}
       {The five ways of reassociating $((A\otimes B)\otimes C)\otimes D$ into
        $A\otimes(B\otimes(C\otimes D))$ using $\alpha$ all agree.}

\qaitem{What does the triangle axiom say?}
       {$(\rho_A \otimes \operatorname{id}_B) = (\operatorname{id}_A \otimes \lambda_B) \circ
        \alpha_{A,I,B}$ as morphisms $(A\otimes I)\otimes B \to A\otimes B$.}

\qaitem{What is a strict monoidal category?}
       {One where $\alpha$, $\lambda$, and $\rho$ are all identity morphisms;
        the tensor product is strictly associative and unital.}

\qaitem{What does Mac Lane's coherence theorem say?}
       {Every monoidal category is monoidally equivalent to a strict one;
        any diagram of associators and unitors between the same two
        parenthesizations commutes.}

\qaitem{What does coherence allow us to do in practice?}
       {Write $A \otimes B \otimes C$ without parentheses and freely move between
        parenthesizations without tracking the specific isomorphisms.}

\qaitem{What is a symmetric monoidal category?}
       {A monoidal category with a natural isomorphism $\sigma_{A,B} : A\otimes B \cong B\otimes A$
        satisfying $\sigma_{B,A} \circ \sigma_{A,B} = \operatorname{id}$ and hexagon axioms.}

\qaitem{What is a closed monoidal category?}
       {A monoidal category where $(-)\otimes B$ has a right adjoint $[B,-]$ (internal hom)
        for every $B$: $\mathcal{V}(A\otimes B, C) \cong \mathcal{V}(A, [B,C])$.}

\qaitem{What is the internal hom in $\mathbf{Set}$?}
       {$[B, C] = C^B$, the set of functions $B \to C$. The adjunction is the
        currying bijection $\operatorname{Hom}(A\times B, C) \cong \operatorname{Hom}(A, C^B)$.}

\qaitem{What is a $\mathcal{V}$-enriched category?}
       {Like a category, but hom-sets $\mathcal{C}(a,b)$ are objects of $\mathcal{V}$;
        composition is a morphism $\mathcal{C}(b,c)\otimes\mathcal{C}(a,b) \to \mathcal{C}(a,c)$;
        identity is a morphism $I \to \mathcal{C}(a,a)$.}

\qaitem{What is a $\mathbf{Set}$-enriched category?}
       {An ordinary category. Hom-sets are sets; composition is a function.}

\qaitem{What is an $\mathbf{Ab}$-enriched category?}
       {A preadditive category: hom-sets are abelian groups and composition
        is bilinear (a $\mathbb{Z}$-bilinear map).}

\qaitem{What is a $\mathbf{Cat}$-enriched category?}
       {A 2-category: hom-objects are categories; objects are 0-cells,
        morphisms are 1-cells, and morphisms between morphisms are 2-cells.}

\qaitem{What is a $([0,\infty], \geq, +, 0)$-enriched category?}
       {A (generalized) metric space: objects are points, $\mathcal{C}(a,b)$ is
        the distance from $a$ to $b$, satisfying $d(a,a) = 0$ and the triangle
        inequality.}

\qaitem{What is a $\mathcal{V}$-functor?}
       {A function on objects plus, for each pair $a,b$, a morphism
        $F_{a,b} : \mathcal{C}(a,b) \to \mathcal{D}(Fa,Fb)$ in $\mathcal{V}$,
        compatible with composition and identities.}

\qaitem{What is a $\mathcal{V}$-natural transformation?}
       {Assigns to each $a$ a morphism $\alpha_a : I \to \mathcal{D}(Fa, Ga)$
        in $\mathcal{V}$, satisfying $\mathcal{V}$-naturality.}

\qaitem{What does a $\mathcal{V}$-natural transformation reduce to when $\mathcal{V} = \mathbf{Set}$?}
       {An ordinary natural transformation: a morphism $I = \{*\} \to \mathcal{D}(Fa,Ga)$
        is just an element of the hom-set.}

\qaitem{What is self-enrichment of a closed monoidal category?}
       {A closed monoidal $\mathcal{V}$ enriches over itself with hom-objects
        $[a,b] \in \mathcal{V}$; composition given by the closed structure.}

\qaitem{What monoidal category is used to enrich toward 2-categories?}
       {$\mathbf{Cat}$ with product $\times$ and unit the terminal category $\mathbf{1}$.}

\qaitem{What is a monoidal functor?}
       {A functor $F : \mathcal{V} \to \mathcal{W}$ between monoidal categories
        with natural morphisms $F(A)\otimes F(B) \to F(A\otimes B)$ and
        $I_\mathcal{W} \to F(I_\mathcal{V})$, compatible with associators and unitors.}

\qaitem{What is the tensor product of abelian groups?}
       {$A \otimes_\mathbb{Z} B$ is the abelian group generated by symbols $a\otimes b$
        subject to bilinearity. It is the unit for the closed structure with
        $[A,B] = \operatorname{Hom}_\mathbb{Z}(A,B)$.}

\qaitem{Why does the monoidal category need to be closed to allow self-enrichment?}
       {The identity element of a $\mathcal{V}$-category requires a map
        $I \to \mathcal{C}(a,a)$; the composition requires a map
        $\mathcal{C}(b,c)\otimes\mathcal{C}(a,b) \to \mathcal{C}(a,c)$.
        In the self-enriched case, both are constructed from the closed structure.}

\qaitem{What is the $\mathcal{V}$-Yoneda lemma?}
       {For a $\mathcal{V}$-category $\mathcal{C}$ and $\mathcal{V}$-functor
        $F : \mathcal{C} \to \mathcal{V}$: $\mathcal{V}\text{-}\mathbf{Nat}(\mathcal{C}(a,-), F) \cong F(a)$
        as objects of $\mathcal{V}$, naturally in $a$ and $F$.}

\qaitem{What is a weighted limit in a $\mathcal{V}$-category?}
       {A limit weighted by a $\mathcal{V}$-functor $W : \mathcal{J} \to \mathcal{V}$;
        generalizes ordinary limits (where $W = \Delta I$) and allows enriched
        universal properties.}

\qaitem{How does enrichment help express that $\operatorname{Hom}(A,B)$ for
        abelian groups is itself an abelian group?}
       {$\mathbf{Ab}$-enrichment of $\mathbf{Ab}$ makes each hom-set $\operatorname{Hom}(A,B)$
        an abelian group, with composition bilinear. This is the self-enriched structure.}

\qaitem{What is a monoid in a monoidal category?}
       {An object $M \in \mathcal{V}$ with morphisms $\mu : M \otimes M \to M$
        (multiplication) and $\eta : I \to M$ (unit), satisfying associativity
        and unit laws. Monoids in $\mathbf{Set}$ are ordinary monoids; in $\mathbf{Ab}$
        they are rings.}

\qaitem{What is a comonoid in a monoidal category?}
       {An object $C$ with comultiplication $\delta : C \to C \otimes C$ and
        counit $\varepsilon : C \to I$, satisfying dual laws. Comonoids in
        $\mathbf{Set}$ with cartesian product are just sets (every set is a
        comonoid via the diagonal).}

\qaitem{What is a braided monoidal category?}
       {A monoidal category with a natural isomorphism $\gamma_{A,B} : A\otimes B \to B\otimes A$
        satisfying hexagon axioms but \emph{not} required to satisfy $\gamma_{B,A} \circ \gamma_{A,B} = \operatorname{id}$.
        Symmetric monoidal is the special case where it does.}
\end{qa}
\end{exercisesbox}

% ─────────────────────────────────────────────────────────────────────────────
\begin{summarybox}
A \textbf{monoidal category} $(\mathcal{V}, \otimes, I)$ is a category with a
tensor product and unit object, with associativity and unitality natural
isomorphisms satisfying the pentagon and triangle axioms. Mac Lane's
\textbf{coherence theorem} says any monoidal category is equivalent to a strict
one, so associativity can be ignored in practice.

A monoidal category is \textbf{symmetric} if $A \otimes B \cong B \otimes A$
naturally, and \textbf{closed} if $(-) \otimes B$ has a right adjoint
$[B, -]$ (internal hom).

A \textbf{$\mathcal{V}$-enriched category} replaces hom-sets with objects
of $\mathcal{V}$, with composition as a morphism in $\mathcal{V}$. $\mathbf{Set}$-categories
are ordinary categories; $\mathbf{Ab}$-categories are preadditive; $\mathbf{Cat}$-categories
are 2-categories; $([0,\infty], +, 0)$-categories are metric spaces. Closed
monoidal categories can enrich over themselves.
\end{summarybox}

% ─────────────────────────────────────────────────────────────────────────────
\section{Formal Restatement}

\begin{formalrestatement}
\textbf{Monoidal category.}
$(\mathcal{V}, \otimes, I, \alpha, \lambda, \rho)$ where
$\otimes : \mathcal{V} \times \mathcal{V} \to \mathcal{V}$,
$I \in \mathcal{V}$, and natural isomorphisms
$\alpha_{A,B,C} : (A\!\otimes\! B)\!\otimes\! C \cong A\!\otimes\!(B\!\otimes\! C)$,
$\lambda_A : I\!\otimes\! A \cong A$,
$\rho_A : A\!\otimes\! I \cong A$
satisfying the pentagon and triangle coherence conditions.

\textbf{Symmetric.} Additional natural isomorphism
$\sigma_{A,B} : A\otimes B \xrightarrow{\sim} B\otimes A$ with
$\sigma_{B,A}\circ\sigma_{A,B} = \operatorname{id}$ and hexagon axioms.

\textbf{Closed.} $(-)\otimes B \dashv [B,-]$ for all $B$; internal hom
$[B,C] \in \mathcal{V}$ with $\mathcal{V}(A\otimes B, C) \cong \mathcal{V}(A,[B,C])$.

\textbf{$\mathcal{V}$-category.} Objects $\operatorname{ob}(\mathcal{C})$; for
each $a,b$: hom-object $\mathcal{C}(a,b)\in\mathcal{V}$; composition
$\circ_{a,b,c} : \mathcal{C}(b,c)\otimes\mathcal{C}(a,b) \to \mathcal{C}(a,c)$;
identity $j_a : I \to \mathcal{C}(a,a)$; satisfying associativity and unit laws
as diagrams in $\mathcal{V}$.

\textbf{$\mathcal{V}$-functor $F : \mathcal{C} \to \mathcal{D}$.}
Function on objects plus $F_{a,b} : \mathcal{C}(a,b) \to \mathcal{D}(Fa,Fb)$
in $\mathcal{V}$, compatible with $\circ$ and $j$.

\textbf{$\mathcal{V}$-natural transformation $\alpha : F \Rightarrow G$.}
Components $\alpha_a : I \to \mathcal{D}(Fa,Ga)$ in $\mathcal{V}$
satisfying: $\circ(G_{a,b}\otimes\alpha_a) = \circ(\alpha_b\otimes F_{a,b})$
as morphisms $\mathcal{C}(a,b)\otimes I \to \mathcal{D}(Fa,Gb)$.

\textbf{$\mathcal{V}$-Yoneda.}
$[\mathcal{C},\mathcal{V}]_\mathcal{V}(\mathcal{C}(a,-),\,F) \cong F(a)$ in $\mathcal{V}$,
naturally in $a \in \mathcal{C}$ and $F \in [\mathcal{C},\mathcal{V}]_\mathcal{V}$.
\end{formalrestatement}
